%! Fundamental Theorem of Algebra and Complex Zeros — Unit 4, Lesson 4.5 — Lesson Plan
\documentclass[10pt]{article}
\usepackage{algebra2-article}
\usepackage{algebra2-boxes}
\usepackage{graphicx}   % -article does NOT load graphicx; needed for thumbnails/screenshots
\graphicspath{{images/}}
\newcommand{\TallMath}[1]{$\displaystyle #1\rule[-1.4em]{0pt}{3.2em}$}
\newcommand{\CourseName}{Algebra 2: Shepherd}
\newcommand{\SchoolYear}{2026--2027}
\newcommand{\MeetingLength}{55 minutes}
\newcommand{\UnitNumberName}{Unit 4: Polynomial Functions \quad}
\newcommand{\LessonNumberName}{Lesson 4.5: Fundamental Theorem of Algebra and Complex Zeros}

\begin{document}
\begin{center}
    {\Large\bfseries \CourseName: \SchoolYear \\
    {\small\bfseries \UnitNumberName \LessonNumberName}}
\end{center}

\begin{tcolorbox}[colback=forestbg, colframe=forest, boxrule=0.9pt, arc=2mm,
    left=2mm, right=2mm, top=1.5mm, bottom=1.5mm]
    \textbf{Primary Objective:} Students can say --- and use --- how many solutions a polynomial
    equation has \emph{before} they solve it. A polynomial of degree $n$ has \textbf{exactly $n$
    complex solutions} when they are counted with multiplicity (the Fundamental Theorem of Algebra),
    and when the coefficients are real the imaginary ones arrive in \textbf{conjugate pairs}. Students
    determine the \emph{number and type} of solutions from a degree and a graph, finish Lesson~4.4's
    workflow \emph{over the complex numbers} instead of stopping at ``no real solutions,'' run the
    process \textbf{backward} to build a polynomial in factored and standard form from a list of
    zeros, and verify their answers by counting against the degree and checking against the graph.
    \\[2pt]
    \textbf{Standards (2023 VA SOL):} \textbf{A2.EI.6b} --- determine the number and type of solutions
    (real or imaginary) of a polynomial equation of degree three or higher --- and \textbf{A2.EI.6a}
    --- determine a factored form of such an equation given its zeros or the $x$-intercepts of its
    graph. \textbf{A2.EI.6c} is now completed \emph{over the set of complex numbers}, and
    \textbf{A2.EI.6d} (verify algebraically and graphically, explain the method) is the counting
    check. Built on \textbf{A2.EO.4a/b/c} (the imaginary unit and $a+bi$ arithmetic, Lesson~3.4),
    \textbf{A2.EI.2b} and the discriminant (Lesson~3.5), \textbf{A2.EO.3b/c} (factoring and division,
    Lessons~4.2--4.3), and \textbf{A2.F.2c} (multiplicity and $x$-intercepts, Lesson~4.0).
\end{tcolorbox}

\renewcommand{\arraystretch}{1.4}
\begin{skillbox}[Priority Ideas \& Skills]{goldbox}
\begin{tabularx}{\linewidth}{@{}X|X@{}}
\textbf{Priority Ideas \& Skills} & \textbf{Key Understandings (the ``why'')} \\[2pt]
\hline
\vspace{-6pt}
\begin{itemize}[leftmargin=1.1em, nosep, after=\vspace{2pt}]
  \item State the \textbf{Fundamental Theorem of Algebra} and its counting corollary: degree $n$ $\Rightarrow$ \emph{exactly} $n$ complex zeros, counted with \textbf{multiplicity}.
  \item Write a polynomial in \textbf{completely factored form over $\mathbb{C}$}: $a(x-r_1)(x-r_2)\cdots(x-r_n)$.
  \item Apply the \textbf{Complex Conjugate Root Theorem}: real coefficients $\Rightarrow$ $a+bi$ is a zero exactly when $a-bi$ is.
  \item Determine the \textbf{number and type} of solutions from the degree plus the $x$-intercepts of a pre-drawn graph (A2.EI.6b).
  \item \textbf{Solve over $\mathbb{C}$}: run Lesson~4.4's workflow, then finish the depressed quadratic with square roots or the quadratic formula even when the discriminant is negative.
  \item \textbf{Build} a polynomial in factored and standard form from a list of zeros, supplying missing conjugates (A2.EI.6a).
  \item \textbf{Verify} by counting against the degree and by checking which zeros the graph can and cannot show.
\end{itemize}
&
\vspace{-6pt}
\begin{itemize}[leftmargin=1.1em, nosep, after=\vspace{2pt}]
  \item Nothing was ever missing. ``No real solutions'' was a statement about $\mathbb{R}$, not about the polynomial.
  \item The degree is a \emph{promise}. It tells you how many answers to hunt for, so ``how do I know I'm done?'' has an arithmetic answer.
  \item Imaginary zeros come in pairs because the conjugate of a real number is itself --- pairing is what keeps the coefficients real.
  \item Therefore an \emph{odd}-degree polynomial with real coefficients must have at least one real zero --- exactly what Lesson~4.0's opposite end-behavior arms forced graphically.
  \item An $x$-intercept is a \emph{real} zero. A graph is an honest but incomplete witness: it can never show you an imaginary zero.
  \item Zeros $\rightarrow$ factors $\rightarrow$ polynomial is the same road as Lesson~4.4, driven in reverse.
\end{itemize}
\end{tabularx}
\end{skillbox}

\begin{skillbox}[{Vocabulary, Concepts, \& Theorems}]{sky}
\renewcommand{\arraystretch}{1.3}
\begin{tabularx}{\linewidth}{@{}l X@{}}
\textbf{Fundamental Theorem of Algebra} & \TallMath{\text{Every polynomial of degree } n \ge 1 \text{ with complex coefficients has at least one complex zero.}} \\
\textbf{Counting corollary} & Degree $n$ $\Rightarrow$ exactly $n$ complex zeros, counted with multiplicity; $f(x)=a(x-r_1)\cdots(x-r_n)$. \\
\textbf{Multiplicity} & How many times a factor repeats; a double zero uses up \emph{two} of the $n$ slots (Lesson 4.0). \\
\textbf{Complex Conjugate Root Theorem} & If the coefficients are \emph{real} and $a+bi$ is a zero, then $a-bi$ is a zero too. \\
\textbf{Complex conjugate pair} & $a+bi$ and $a-bi$; their factors multiply to the \emph{real} quadratic $x^2-2ax+(a^2+b^2)$. \\
\textbf{Number and type} & How many solutions are real and how many are imaginary --- answered before solving (A2.EI.6b). \\
\textbf{Real zero vs.\ $x$-intercept} & The same object. Imaginary zeros are \emph{not} on the graph, which is why intercepts can be fewer than the degree. \\
\textbf{Factored form over $\mathbb{C}$} & Every factor linear. Over $\mathbb{R}$ an irreducible quadratic like $x^2+4$ is as far as you can go. \\
\end{tabularx}
\end{skillbox}

\begin{fixedskillbox}[Activate Prior Knowledge \& Spiral Review]{forestbg}
    The Warm-Up rebuilds the three tools this lesson needs and then hands students the answer to
    yesterday's cliffhanger. \textbf{(1)} \emph{Complex arithmetic from Lesson~3.4}: $i^2=-1$ and the
    conjugate product $(x-3i)(x+3i)=x^2+9$ --- with the observation that the $i$'s vanish and the
    result is \emph{real}. \textbf{(2)} \emph{Solving from Lessons~3.3 and 3.5}: $x^2+4=0$ by square
    roots ($x=\pm 2i$) and $x^2-2x+5=0$ by the quadratic formula ($D=-16$, $x=1\pm 2i$) --- both
    answers written down as a \emph{pair}. \textbf{(3)} \emph{Finishing Lesson~4.4's Tier E problem}:
    $x^3-x^2+4x-4=(x-1)(x^2+4)$, whose depressed polynomial stalled yesterday, now solved completely
    to $1$, $2i$, $-2i$. \textbf{(4)} The count: a cubic that gave up one solution yesterday has
    \emph{three} today. Debrief items 1 and 3 together --- the conjugates in item~1 are not a
    coincidence --- and do \emph{not} name the two theorems yet.
\end{fixedskillbox}

\begin{skillbox}[Hook]{forestbg}
    Put $f(x) = x^3-x^2+4x-4$ on the board with its graph (pre-drawn in the notes). It crosses the
    $x$-axis \textbf{exactly once}, at $x=1$. Ask what the degree is, then ask the question the whole
    lesson answers: \textbf{``A cubic is supposed to have three solutions. This graph shows one. Where
    are the other two?''} Let the class argue. Then point at the Warm-Up: they already found them ---
    $2i$ and $-2i$. \textbf{``So they exist. They are just not on the graph, because a graph only
    shows you \emph{real} zeros.''} Reframe every ``no solution'' moment of the year --- $x^2=-4$ in
    Lesson~3.3, the negative discriminant in 3.5, yesterday's stalled depressed polynomial --- as the
    same sentence: \emph{no real solutions}, which was never the same as \emph{no solutions}. Promise
    the payoff: by the end of the period they will know how many answers a polynomial equation has
    before they start solving it, and they will be able to build a polynomial to order from a list of
    zeros.
\end{skillbox}

\begin{skillbox}[Lesson]{forestbg}
\begin{multicols}{2}
\textbf{1. The theorem, and the corollary that does the work.} State the Fundamental Theorem of
Algebra --- every polynomial of degree $n\ge1$ has at least one complex zero --- and then get quickly
to the version students will actually use: apply it, divide out the zero, apply it again, and the
process runs exactly $n$ times, so a degree-$n$ polynomial factors completely as
$a(x-r_1)\cdots(x-r_n)$ and has \emph{exactly $n$} zeros counted with multiplicity. Emphasize the
word \emph{exactly}: this is the first tool of the year that says how many answers there are
\emph{before} any are found.

\textbf{2. Counting with multiplicity.} Return to the unit anchor $g(x)=x^3-3x+2=(x-1)^2(x+2)$, whose
graph students read in Lesson~4.0. Two \emph{distinct} zeros, three zeros \emph{counted with
multiplicity}: the double zero at $1$ (where the graph touches) fills two of the three slots. Say the
sentence to memorize: \textbf{``Distinct zeros can be fewer than $n$. Zeros counted with multiplicity
are never fewer, and never more.''}

\textbf{3. Imaginary zeros travel in pairs.} Look back at Warm-Up item~2: $1+2i$ and $1-2i$; and at
item~3: $2i$ and $-2i$. State the Complex Conjugate Root Theorem for \emph{real} coefficients, and
give the one-line reason --- the conjugate factors multiply to $x^2-2ax+(a^2+b^2)$, which is real, so
the pair is exactly what lets a real polynomial hide an imaginary zero. Then draw the two
consequences: the number of imaginary zeros is always \textbf{even}, and therefore an \textbf{odd}
degree forces at least one real zero. Tie that to Lesson~4.0: an odd-degree graph has arms pointing
opposite directions, so it \emph{has} to cross.

\columnbreak

\textbf{4. Number and type, from a graph (A2.EI.6b).} The section that most directly matches the
standard, and it is pure graph reading. Degree gives the total; $x$-intercepts give the real ones (in
multiplicity); subtract for the imaginary ones. Work the table in the notes aloud --- a quartic
showing two intercepts has two imaginary zeros; a cubic showing one has two. Ask the killer question
at least twice: \textbf{``Could a cubic with real coefficients show \emph{zero} $x$-intercepts?''}
(No --- that would need three imaginary zeros, and they come in pairs.)

\textbf{5. Solve over $\mathbb{C}$ --- the workflow, finished.} Nothing new procedurally: list, test,
divide, finish (Lesson~4.4). The change is the last step. When the depressed quadratic has a negative
discriminant you no longer write ``no solution'' --- you write the conjugate pair. Do
$x^3-5x^2+9x-5=0$ ($x=1$, then $x^2-4x+5$, then $2\pm i$) so students see a pair whose real part is
\emph{not} zero, and $x^4-5x^2-36=0$ by quadratic form (Lesson~4.2) for $\pm3, \pm2i$.

\textbf{6. Run it backward (A2.EI.6a).} Zeros $\rightarrow$ factors $\rightarrow$ multiply. Build
$(x-2)(x^2+9)=x^3-2x^2+9x-18$ from the zeros $2$ and $3i$, insisting on \emph{why} $-3i$ had to be
invited. Then $(x+1)^2(x-4)$ from a repeated zero, and close with the question that tests whether the
theorem is understood: \emph{what is the smallest possible degree of a real polynomial with zeros $5$
and $1-i$?}
\end{multicols}
\end{skillbox}

\begin{skillbox}[Active Monitoring]{redbox}
Circulate during the notes practice and Tier work. Watch for and correct:
\begin{itemize}[leftmargin=1.2em, nosep]
  \item \textbf{a lone imaginary zero} --- a real polynomial written with $3i$ but not $-3i$. This is the signature error of the lesson; ask them to multiply out and find the $i$ left in the coefficients;
  \item \textbf{counting distinct zeros instead of counting with multiplicity} --- ``$(x-1)^2(x+2)$ has two zeros, so the theorem is wrong'';
  \item \textbf{looking for imaginary zeros on the graph} --- ``I can't find the other two intercepts.'' They are not there and never will be;
  \item \textbf{stopping at ``no real solutions''} on a depressed quadratic --- the exact habit this lesson replaces;
  \item \textbf{sign slips in the conjugate product} --- $(x-3i)(x+3i)$ becoming $x^2-9$; make them show the $-9i^2$ step;
  \item \textbf{$\pm$ dropped when taking a square root} of a negative number ($x^2=-4 \Rightarrow x=2i$ only);
  \item \textbf{forgetting the ``real coefficients'' hypothesis} when they quote the conjugate theorem;
  \item building a polynomial and \textbf{never multiplying out} --- A2.EI.6a asks for a factored form, but the check that it is right lives in standard form.
\end{itemize}
Cold-call prompts: ``What is the degree, and how many solutions does that promise?'' ``You wrote $2i$
--- who has to come with it, and why?'' ``This quartic's graph crosses twice. How many imaginary
zeros does it have?'' ``Can an odd-degree polynomial have \emph{no} real zeros? Convince me two ways
--- pairing, and end behavior.'' ``Your depressed polynomial is $x^2+9$. Is that the end of the
problem or the middle?''
\end{skillbox}

\begin{skillbox}[Group Work \& Differentiation]{redbox}
\textbf{Nothing Is Missing} activity (tiered; mirrors the notes).
\begin{multicols}{3}
\textbf{Tier R --- Remediate.} Counting first, solving second. Students read four already-factored
polynomials and report degree, distinct zeros, and zeros counted with multiplicity --- including one
with an $x^2+4$ factor, so a conjugate pair gets counted without any solving. Then two quadratics over
$\mathbb{C}$ ($x^2+25=0$ and $x^2-6x+13=0$) rebuild the two finishing moves, and a guided solve of
$x^3-2x^2+9x-18=0$ (candidate given, synthetic grid pre-drawn) produces $2$, $\pm3i$ --- the same
polynomial the notes \emph{built} from those zeros, met from the other direction. Closes with
building $x(x+3)(x-2)$ from three real zeros.

\columnbreak
\textbf{Tier A --- Approaching.} A full solve over $\mathbb{C}$ with no candidate handed over,
$x^3-3x^2+4x-12=0$ (zeros $3$, $\pm2i$); a \textbf{graph-driven} number-and-type problem on a
pre-drawn quartic with exactly two $x$-intercepts, where students must report two real and two
imaginary zeros \emph{without solving anything}; building a cubic from the zeros $5$ and $1-i$
(supplying the conjugate, then multiplying to $x^3-7x^2+12x-10$); and error analysis on the two
signature mistakes --- an unpaired imaginary zero, and hunting for imaginary zeros among the
$x$-intercepts.

\columnbreak
\textbf{Tier E --- Extension.} A guided proof that the conjugate pair always produces the \emph{real}
quadratic $x^2-2ax+(a^2+b^2)$, and the argument it powers: why an odd-degree real polynomial must have
a real zero, cross-checked against Lesson~4.0's end behavior. Then $x^4+13x^2+36=0$, whose four zeros
$\pm2i,\pm3i$ are \emph{all} imaginary, with a pre-drawn graph that never touches the $x$-axis;
building a quartic from the zeros $2i$ and $1+i$ (both need partners, so the degree is forced); and a
sharp counterexample --- $f(x)=x-i$ has the zero $i$ and not $-i$, which is exactly why the conjugate
theorem says \emph{real} coefficients.
\end{multicols}
\end{skillbox}

\begin{skillbox}[Individual Work \& Assessment]{redbox}
\textbf{Exit Ticket} (independent, no notes): \textbf{(1)} solve $x^3+x^2+4x+4=0$ over $\mathbb{C}$
--- factor by grouping (Lesson~4.2) to $(x+1)(x^2+4)$, report $-1$, $2i$, $-2i$, and then state the
\emph{number and type} explicitly (one real, two imaginary) and how the degree confirms the list is
complete; \textbf{(2)} write a polynomial of \emph{least} degree with real coefficients whose zeros
include $-2$ and $4i$ --- the item that fails immediately if the conjugate is not invited, answer
$(x+2)(x^2+16)=x^3+2x^2+16x+32$; and \textbf{(3)} an \textbf{SOL-style multiple-choice} item asking
which list \emph{cannot} be the complete set of zeros of a degree-4 polynomial with real coefficients,
with an unpaired $3i$ as the answer and three legitimate lists (including one using multiplicity) as
distractors. Sort into ``got it / stopped at no-real-solutions / conjugate not supplied / counted
wrong.'' The conjugate pile is the one to fix on the spot, because Lesson~4.6 asks students to move
between a graph, a factored form, and a degree constantly, and a missing conjugate breaks that
triangle.
\end{skillbox}

\begin{skillbox}[Reinforcement \& Extension]{goldbox}
\textbf{Homework:} counting practice from factored forms (including a repeated \emph{imaginary} zero);
number-and-type from a degree plus a described set of $x$-intercepts; three solves over $\mathbb{C}$
of escalating difficulty --- $x^3+2x^2+4x+8=0$ (the polynomial Lesson~4.3's homework produced as a
quotient of $x^4-16$, now finished to $-2,\pm2i$), $x^3-7x^2+17x-15=0$ (zeros $3$, $2\pm i$), and the
quartic $x^4-16=0$ itself, whose four zeros $\pm2,\pm2i$ close a loop opened two lessons ago; two
build-from-zeros items, one with a repeated zero and one requiring a supplied conjugate; a
degree-5 item where three zeros and a multiplicity are given and the fifth must be deduced; and two
written items --- why an odd-degree polynomial always has a real zero, and why a graph is an
incomplete witness. \textbf{Extension:} prove the conjugate factors always give real coefficients;
find the smallest degree of a real polynomial with prescribed zeros; and a design task --- write a
quartic with exactly two $x$-intercepts, then one with none, and justify both with the counting rules.
\textbf{Preview:} Lesson~4.6, graphing polynomial functions and modeling --- everything the unit has
built (end behavior, multiplicity, zeros, and now the complete count) becomes one procedure for
producing and reading a polynomial's graph.
\end{skillbox}
\end{document}
